\ifdefined\BookMaster
\def\BookEnd{}
\else
\documentclass[a4paper,12pt,fleqn,openany,twoside,fontset=windows]{ctexbook}

\usepackage[fleqn]{amsmath}
\usepackage{amssymb,amsthm,mathrsfs}
\usepackage{geometry}
\usepackage[dvipsnames]{xcolor}
\usepackage{graphicx}
\usepackage{tikz}
\usetikzlibrary{calc}
\usepackage[most]{tcolorbox}
\usepackage{fancyhdr}
\usepackage{titlesec}
\usepackage{tocloft}
\usepackage{enumitem}
\usepackage{microtype}
\usepackage{pifont}
\usepackage{hyperref}
\usepackage{romannum}
\usepackage{mathtools}

\definecolor{bookblue}{HTML}{264653}
\definecolor{bookteal}{HTML}{4F7C82}
\definecolor{booklight}{HTML}{EDF4F3}
\definecolor{bookgray}{HTML}{5E686B}

\geometry{
    inner=2.7cm,
    outer=2.5cm,
    top=2.7cm,
    bottom=2.7cm,
    headheight=18pt,
    headsep=18pt,
    footskip=25pt
}
\setlength{\mathindent}{0pt}
\setlength{\parindent}{5pt}
\setlength{\parskip}{3pt plus 1pt minus 1pt}
\linespread{1.25}
\allowdisplaybreaks[3]
\AtBeginDocument{%
    \setlength{\parindent}{5pt}
    \setlength{\abovedisplayskip}{8pt plus 2pt minus 2pt}
    \setlength{\belowdisplayskip}{8pt plus 2pt minus 2pt}
    \setlength{\abovedisplayshortskip}{5pt plus 1pt minus 1pt}
    \setlength{\belowdisplayshortskip}{5pt plus 1pt minus 1pt}
}

\hypersetup{
    unicode=true,
    colorlinks=true,
    linkcolor=bookblue,
    citecolor=bookblue,
    urlcolor=bookteal,
    bookmarksopen=true,
    pdfauthor={Chen},
    pdftitle={数学分析习题整理}
}

\renewcommand{\chaptermark}[1]{%
    \edef\@tmp{第\thechapter 章\quad #1}%
    \expandafter\markboth\expandafter{\@tmp}{}%
}
\fancypagestyle{main}{%
    \fancyhf{}
    \fancyhead[LE]{\small\color{bookgray}\nouppercase{\leftmark}}
    \fancyhead[RO]{\small\color{bookgray}数学分析习题整理}
    \fancyfoot[C]{\color{bookblue}\thepage}
    \fancyfoot[R]{\small\bfseries Chen\;\; github.com/mathlover-1/math-analysis}
    \renewcommand{\headrulewidth}{0.4pt}
    \renewcommand{\footrulewidth}{0pt}
}
\fancypagestyle{plain}{%
    \fancyhf{}
    \fancyfoot[C]{\color{bookblue}\thepage}
    \fancyfoot[R]{\small\bfseries Chen\;\; github.com/mathlover-1/math-analysis}
    \renewcommand{\headrulewidth}{0pt}
    \renewcommand{\footrulewidth}{0pt}
}
\pagestyle{main}

\titleformat{\chapter}[display]
  {\normalfont\sffamily\bfseries\color{bookblue}}
  {\raggedleft\fontsize{46}{50}\selectfont\thechapter}
  {-18pt}
  {\titlerule[1.2pt]\vspace{10pt}\filright\Huge}
\titlespacing*{\chapter}{0pt}{-15pt}{36pt}

\renewcommand{\contentsname}{目\quad 录}
\renewcommand{\cfttoctitlefont}{\hfill\sffamily\bfseries\Huge\color{bookblue}}
\renewcommand{\cftaftertoctitle}{\hfill}
\setlength{\cftbeforetoctitleskip}{5pt}
\setlength{\cftaftertoctitleskip}{28pt}
\setlength{\cftbeforechapskip}{12pt}
\renewcommand{\cftchapfont}{\sffamily\bfseries\large\color{bookblue}}
\renewcommand{\cftchappagefont}{\sffamily\bfseries\large\color{bookblue}}
\renewcommand{\cftchapleader}{\color{bookteal}\cftdotfill{2.5}}

\newtcolorbox{ti}{
    enhanced,
    breakable,
    colback=booklight!55!white,
    colframe=bookteal!35!white,
    boxrule=0.45pt,
    boxsep=0pt,
    left=10pt,
    right=10pt,
    top=9pt,
    bottom=9pt,
    before skip=14pt,
    after skip=14pt,
    arc=2pt,
    outer arc=2pt,
    before upper={\parindent=0pt}
}

\newenvironment{booknotebody}{%
    \begin{center}
    \begin{minipage}{0.84\textwidth}
    \songti\zihao{-4}\linespread{1.45}\selectfont
    \setlength{\parindent}{2\ccwd}
    \setlength{\parskip}{0.75em}
}{%
    \end{minipage}
    \end{center}
}

\newenvironment{booknotelongbody}{%
    \begingroup
    \songti\zihao{-4}\linespread{1.45}\selectfont
    \setlength{\parindent}{2\ccwd}
    \setlength{\parskip}{0.75em}
    \begin{list}{}{%
        \setlength{\leftmargin}{0.08\textwidth}
        \setlength{\rightmargin}{0.08\textwidth}
        \setlength{\topsep}{0pt}
        \setlength{\partopsep}{0pt}
        \setlength{\parsep}{0pt}
        \setlength{\itemsep}{0pt}
    }
    \item\relax
}{%
    \end{list}
    \endgroup
}

\fancypagestyle{plainnowm}{%
    \fancyhf{}
    \fancyfoot[C]{\color{bookblue}\thepage}
    \renewcommand{\headrulewidth}{0pt}
    \renewcommand{\footrulewidth}{0pt}
}

\newcommand{\booknotetitle}[2][plain]{%
    \clearpage
    \phantomsection
    \addcontentsline{toc}{chapter}{#2}
    \markboth{#2}{}
    \thispagestyle{#1}
    \vspace*{1.2cm}
    \begin{center}
        {\sffamily\heiti\bfseries\fontsize{26}{34}\selectfont\color{bookblue}#2\par}
        \vspace{0.8em}
        {\color{bookteal}\rule{4.5cm}{0.8pt}\par}
    \end{center}
    \vspace{0.6cm}
}

\renewcommand{\proofname}{证明}
\renewcommand{\qedsymbol}{\color{bookteal}\ensuremath{\blacksquare}}

\title{数学分析习题整理}
\author{Chen}
\date{}

\makeatletter
% ==================== 旧版封面/封底设计备份 ====================
% 2026-08 重设计,旧代码用 \iffalse...\fi 封存,新版见下方。
\iffalse
\renewcommand{\maketitle}{%
    \begin{titlepage}
        \thispagestyle{empty}
        \setcounter{page}{0}
        \noindent\hspace*{-\parindent}\begin{tikzpicture}[remember picture,overlay]
            \fill[bookblue] (current page.north west) rectangle ([yshift=-4.2cm]current page.north east);
            \fill[booklight] ([yshift=-4.2cm]current page.north west) rectangle (current page.south east);
            \draw[bookteal!55,line width=0.7pt]
                ([xshift=2.4cm,yshift=4.3cm]current page.south west)
                .. controls +(2.5cm,1.8cm) and +(-2.8cm,-1.3cm) ..
                ([xshift=-2.4cm,yshift=7.2cm]current page.south east);
            \draw[bookteal!35,line width=0.7pt]
                ([xshift=2.5cm,yshift=6.2cm]current page.south west)
                .. controls +(3.0cm,-1.2cm) and +(-3.0cm,1.6cm) ..
                ([xshift=-2.3cm,yshift=3.8cm]current page.south east);
        \end{tikzpicture}
        \vspace*{5.6cm}
        \noindent\color{bookblue}\rule{\textwidth}{1.4pt}\par
        \vspace{8pt}
        {\centering\sffamily\bfseries\fontsize{32}{42}\selectfont\color{bookblue}\@title\par}
        \vspace{8pt}
        \noindent\color{bookteal}\rule{\textwidth}{0.6pt}\par
        \vspace{2.6cm}
        {\centering\Large\color{bookgray}\@author\par}
        \vfill
        {\centering\color{bookteal}\large
            $\displaystyle \lim_{n\to\infty} a_n$\qquad
            $\displaystyle \int_a^b f(x)\,\mathrm{d}x$\qquad
            $\displaystyle \nabla f(x,y)$\par
        }
        \vspace*{1.5cm}
    \end{titlepage}
}

\newcommand{\makebackcover}{%
    \begingroup
    \let\cleardoublepage\clearpage
    \begin{titlepage}
        \thispagestyle{empty}
        \noindent\hspace*{-\parindent}\begin{tikzpicture}[remember picture,overlay]
            \fill[bookblue] (current page.north west) rectangle (current page.south east);
            \fill[booklight,rounded corners=7pt]
                ([xshift=2.2cm,yshift=-3.0cm]current page.north west)
                rectangle
                ([xshift=-2.2cm,yshift=3.0cm]current page.south east);
            \fill[bookteal!80]
                (current page.north west) --
                ([xshift=5.3cm]current page.north west) --
                ([xshift=2.1cm,yshift=-4.4cm]current page.north west) --
                ([yshift=-6.2cm]current page.north west) -- cycle;
            \fill[bookteal!55]
                (current page.south east) --
                ([xshift=-5.8cm]current page.south east) --
                ([xshift=-2.4cm,yshift=4.2cm]current page.south east) --
                ([yshift=6.0cm]current page.south east) -- cycle;
            \foreach \radius in {0.75,1.15,1.55,1.95} {
                \draw[bookteal!55,line width=0.55pt]
                    ([xshift=-2.9cm,yshift=-4.4cm]current page.north east)
                    circle[radius=\radius cm];
            }
            \foreach \radius in {0.65,1.05,1.45} {
                \draw[booklight!65,line width=0.55pt]
                    ([xshift=2.4cm,yshift=3.6cm]current page.south west)
                    circle[radius=\radius cm];
            }
            \fill[bookteal!16,rounded corners=3pt]
                ([xshift=3.5cm,yshift=-7.1cm]current page.north west)
                rectangle
                ([xshift=-5.0cm,yshift=7.2cm]current page.south east);
            \fill[bookteal!30,rounded corners=3pt]
                ([xshift=4.4cm,yshift=-8.0cm]current page.north west)
                rectangle
                ([xshift=-4.1cm,yshift=8.1cm]current page.south east);
            \begin{scope}[shift={([xshift=-0.75cm]current page.center)}]
                \clip[rounded corners=3pt] (-5.25,-4.75) rectangle (5.25,4.75);
                \foreach \v in {-3.6,-3.0,...,3.6} {
                    \draw[bookteal!48,line width=0.45pt,opacity=0.72]
                        plot[domain=-3.6:3.6,samples=45,smooth,variable=\u]
                        ({\u+0.55*\v},{0.28*\u-0.35*\v+0.11*(\u*\u-\v*\v)});
                }
                \foreach \u in {-3.6,-3.0,...,3.6} {
                    \draw[bookblue!42,line width=0.42pt,opacity=0.62]
                        plot[domain=-3.6:3.6,samples=45,smooth,variable=\v]
                        ({\u+0.55*\v},{0.28*\u-0.35*\v+0.11*(\u*\u-\v*\v)});
                }
                \draw[bookteal!70,line width=0.75pt,opacity=0.82]
                    plot[domain=-3.6:3.6,samples=55,smooth,variable=\u]
                    ({\u},{0.28*\u+0.11*\u*\u});
                \draw[bookblue!60,line width=0.7pt,opacity=0.75]
                    plot[domain=-3.6:3.6,samples=55,smooth,variable=\v]
                    ({0.55*\v},{-0.35*\v-0.11*\v*\v});
            \end{scope}
            \draw[bookteal!45,line width=0.65pt]
                ([xshift=3.0cm,yshift=-8.2cm]current page.north west)
                .. controls +(2.2cm,2.0cm) and +(-2.0cm,-2.1cm) ..
                ([xshift=-3.0cm,yshift=9.0cm]current page.south east);
            \draw[bookteal!25,line width=0.55pt]
                ([xshift=3.1cm,yshift=8.0cm]current page.south west)
                .. controls +(2.7cm,-1.7cm) and +(-2.6cm,1.8cm) ..
                ([xshift=-3.1cm,yshift=-8.7cm]current page.north east);
            \node[anchor=west,text=bookteal!72,scale=0.78] at
                ([xshift=2.8cm,yshift=-4.0cm]current page.north west) {%
                $\displaystyle e^{\mathrm{i}\pi}+1=0$};
            \node[anchor=east,text=bookblue!62,scale=0.67] at
                ([xshift=-2.7cm,yshift=4.0cm]current page.south east) {%
                $\displaystyle
                f(x)=\sum_{n=0}^{\infty}\frac{f^{(n)}(a)}{n!}(x-a)^n$};
        \end{tikzpicture}
        \mbox{}
    \end{titlepage}
    \endgroup
}
\fi

% ==================== 新版封面(2026-08 重设计) ====================
% 设计要点:顶部深绿纵向渐变+波浪下缘+两层半透明过渡波(解决深浅绿生硬过渡);
% 深色区螺旋线/同心圆母题;标题+英文副题+作者居中;
% 中部 Fourier 部分和逼近方波曲线;底部浅色波纹(无文字)。
\renewcommand{\maketitle}{%
    \begin{titlepage}
        \thispagestyle{empty}
        \setcounter{page}{0}
        \noindent\hspace*{-\parindent}\begin{tikzpicture}[remember picture,overlay]
            % ---------- 底色 ----------
            \fill[booklight] (current page.north west) rectangle (current page.south east);

            % ---------- 顶部深色区:纵向渐变 + 波浪下缘 ----------
            \shade[top color=bookblue, bottom color=bookteal!55!bookblue]
                (current page.north west) -- (current page.north east) --
                ([yshift=-8.6cm]current page.north east)
                .. controls ([xshift=-5.2cm,yshift=-10.6cm]current page.north east)
                         and ([xshift=6.0cm,yshift=-8.0cm]current page.north west) ..
                ([yshift=-10.0cm]current page.north west) -- cycle;

            % ---------- 过渡层 1:半透明波浪 ----------
            \fill[bookteal!60!bookblue,opacity=0.42]
                ([yshift=-8.6cm]current page.north east)
                .. controls ([xshift=-5.2cm,yshift=-10.6cm]current page.north east)
                         and ([xshift=6.0cm,yshift=-8.0cm]current page.north west) ..
                ([yshift=-10.0cm]current page.north west) --
                ([yshift=-12.6cm]current page.north west)
                .. controls ([xshift=6.5cm,yshift=-11.0cm]current page.north west)
                         and ([xshift=-5.5cm,yshift=-13.2cm]current page.north east) ..
                ([yshift=-11.4cm]current page.north east) -- cycle;

            % ---------- 过渡层 2:更浅的波浪 ----------
            \fill[bookteal!35,opacity=0.35]
                ([yshift=-10.6cm]current page.north east)
                .. controls ([xshift=-6.0cm,yshift=-12.4cm]current page.north east)
                         and ([xshift=5.0cm,yshift=-9.8cm]current page.north west) ..
                ([yshift=-11.8cm]current page.north west) --
                ([yshift=-13.8cm]current page.north west)
                .. controls ([xshift=5.5cm,yshift=-12.6cm]current page.north west)
                         and ([xshift=-6.5cm,yshift=-14.6cm]current page.north east) ..
                ([yshift=-12.8cm]current page.north east) -- cycle;

            % ---------- 深色区装饰:右上螺旋线 ----------
            \begin{scope}[shift={([xshift=-3.4cm,yshift=-4.3cm]current page.north east)}]
                \draw[booklight!45!bookteal,opacity=0.55,line width=0.6pt]
                    plot[domain=0:1440,samples=220,smooth,variable=\t]
                    ({0.0028*\t*cos(\t)},{0.0028*\t*sin(\t)});
            \end{scope}

            % ---------- 深色区装饰:左下同心圆 ----------
            \foreach \r in {0.7,1.25,1.8,2.35,2.9}{
                \draw[booklight!35!bookteal,opacity=0.5,line width=0.5pt]
                    ([xshift=0.4cm,yshift=-7.9cm]current page.north west) circle[radius=\r cm];
            }

            % ---------- 主标题 ----------
            \node[anchor=center,text=booklight] at
                ([yshift=-4.5cm]current page.north)
                {\sffamily\bfseries\fontsize{40}{46}\selectfont \@title};

            % ---------- 分隔短线 ----------
            \draw[booklight!65,line width=1.2pt]
                ([xshift=-2.6cm,yshift=-5.75cm]current page.north) --
                ([xshift=2.6cm,yshift=-5.75cm]current page.north);

            % ---------- 英文副题 ----------
            \node[anchor=center,text=booklight!70!bookteal] at
                ([yshift=-6.55cm]current page.north)
                {\itshape\large Mathematical Analysis \textperiodcentered\
                 A Collection of Exercises};

            % ---------- 作者(紧随标题下方) ----------
            \node[anchor=center,text=booklight!88] at
                ([yshift=-8.2cm]current page.north)
                {\sffamily\large \@author\quad 整理};

            % ---------- 中部:Fourier 部分和逼近方波 ----------
            \begin{scope}[shift={([xshift=-5.2cm,yshift=-18.1cm]current page.north)}]
                % 坐标轴
                \draw[bookgray!55,line width=0.5pt,->] (-0.5,0) -- (10.9,0);
                \draw[bookgray!55,line width=0.5pt,->] (0,-1.75) -- (0,1.95);
                % 目标方波(虚线)
                \draw[bookgray!75,line width=0.6pt,dashed]
                    (0,0.92) -- (5,0.92) (5,-0.92) -- (10,-0.92);
                % N=1
                \draw[bookteal!45,line width=0.7pt]
                    plot[domain=0:360,samples=90,smooth,variable=\x]
                    ({\x/36},{1.18*sin(\x)});
                % N=3
                \draw[bookteal!80,line width=0.8pt]
                    plot[domain=0:360,samples=120,smooth,variable=\x]
                    ({\x/36},{1.18*(sin(\x)+sin(3*\x)/3)});
                % N=7
                \draw[bookblue,line width=0.95pt]
                    plot[domain=0:360,samples=150,smooth,variable=\x]
                    ({\x/36},{1.18*(sin(\x)+sin(3*\x)/3+sin(5*\x)/5+sin(7*\x)/7)});
                % 刻度标注
                \node[anchor=north east,text=bookgray!80,scale=0.62] at (-0.12,0) {$0$};
                \node[anchor=north,text=bookgray!80,scale=0.62] at (5,0) {$\pi$};
                \node[anchor=north,text=bookgray!80,scale=0.62] at (10,0) {$2\pi$};
            \end{scope}
            % 图注
            \node[anchor=center,text=bookgray,scale=0.8] at
                ([yshift=-20.5cm]current page.north)
                {$S_N(x)=\displaystyle\sum_{k=1}^{N}\frac{\sin\bigl((2k-1)x\bigr)}{2k-1}$
                 \ 逐步逼近方波};

            % ---------- 底部浅色波纹(无文字,与顶部波浪呼应) ----------
            \fill[bookteal!22,opacity=0.5]
                ([yshift=2.4cm]current page.south west)
                .. controls ([xshift=6.0cm,yshift=3.9cm]current page.south west)
                         and ([xshift=-6.5cm,yshift=1.0cm]current page.south east) ..
                ([yshift=2.9cm]current page.south east) --
                (current page.south east) -- (current page.south west) -- cycle;
            \fill[bookteal!40!bookblue,opacity=0.28]
                ([yshift=1.3cm]current page.south west)
                .. controls ([xshift=5.5cm,yshift=2.7cm]current page.south west)
                         and ([xshift=-5.5cm,yshift=0.3cm]current page.south east) ..
                ([yshift=1.9cm]current page.south east) --
                (current page.south east) -- (current page.south west) -- cycle;
        \end{tikzpicture}
        \mbox{}
    \end{titlepage}
}

% ==================== 新版封底(纯背景,2026-08 重设计) ====================
% 设计要点:对角渐变底+角部半透明叠层;中下部外摆线(玫瑰形)为视觉主体;
% 底部两层半透明白色波浪;无任何文字。
\newcommand{\makebackcover}{%
    \begingroup
    \let\cleardoublepage\clearpage
    \begin{titlepage}
        \thispagestyle{empty}
        \noindent\hspace*{-\parindent}\begin{tikzpicture}[remember picture,overlay]
            % ---------- 底色:对角渐变 ----------
            \shade[top color=bookblue, bottom color=bookteal!48!bookblue,
                   shading angle=-30]
                (current page.north west) rectangle (current page.south east);

            % ---------- 角部半透明叠层 ----------
            \fill[white,opacity=0.05]
                ([xshift=2.0cm,yshift=1.0cm]current page.north east) circle[radius=7.2cm];
            \fill[white,opacity=0.06]
                ([xshift=-0.5cm,yshift=-1.5cm]current page.north east) circle[radius=3.6cm];

            % ---------- 左下同心圆 ----------
            \foreach \r in {1.0,1.75,2.5,3.25,4.0,4.75}{
                \draw[booklight!40,opacity=0.55,line width=0.55pt]
                    ([xshift=-0.6cm,yshift=-0.4cm]current page.south west) circle[radius=\r cm];
            }
            % ---------- 右下螺旋线 ----------
            \begin{scope}[shift={([xshift=-2.9cm,yshift=3.2cm]current page.south east)}]
                \draw[booklight!45!bookteal,opacity=0.5,line width=0.55pt]
                    plot[domain=0:1260,samples=200,smooth,variable=\t]
                    ({0.0036*\t*cos(\t)},{0.0036*\t*sin(\t)});
            \end{scope}
            % ---------- 左上斜线组 ----------
            \foreach \i in {0,1,2}{
                \draw[bookteal!60!booklight,opacity=0.5,line width=0.6pt]
                    ([xshift={1.6cm+\i*0.45cm}]current page.north west) --
                    ([yshift={-1.6cm-\i*0.45cm}]current page.north west);
            }

            % ---------- 中下部视觉主体:外摆线(玫瑰形) ----------
            \begin{scope}[shift={([yshift=-16.6cm]current page.north)}]
                % 外圈
                \draw[booklight!35,opacity=0.6,line width=0.6pt]
                    (0,0) circle[radius=4.9cm];
                \draw[booklight!22,opacity=0.5,line width=0.5pt]
                    (0,0) circle[radius=5.25cm];
                % 主玫瑰线
                \draw[booklight!55!bookteal,opacity=0.8,line width=0.85pt]
                    plot[domain=0:1080,samples=400,smooth,variable=\t]
                    ({0.335*(8*cos(\t)-5*cos(2.6667*\t))},
                     {0.335*(8*sin(\t)-5*sin(2.6667*\t))});
                % 内层玫瑰线
                \draw[booklight!38!bookteal,opacity=0.6,line width=0.6pt]
                    plot[domain=0:1080,samples=400,smooth,variable=\t]
                    ({0.21*(8*cos(\t)-5*cos(2.6667*\t))},
                     {0.21*(8*sin(\t)-5*sin(2.6667*\t))});
            \end{scope}

            % ---------- 底部半透明波浪 ----------
            \fill[white,opacity=0.06]
                ([yshift=4.6cm]current page.south west)
                .. controls ([xshift=6.5cm,yshift=6.4cm]current page.south west)
                         and ([xshift=-6.0cm,yshift=2.6cm]current page.south east) ..
                ([yshift=5.2cm]current page.south east) --
                (current page.south east) -- (current page.south west) -- cycle;
            \fill[booklight,opacity=0.08]
                ([yshift=2.6cm]current page.south west)
                .. controls ([xshift=5.5cm,yshift=4.3cm]current page.south west)
                         and ([xshift=-6.5cm,yshift=1.0cm]current page.south east) ..
                ([yshift=3.3cm]current page.south east) --
                (current page.south east) -- (current page.south west) -- cycle;
        \end{tikzpicture}
        \mbox{}
    \end{titlepage}
    \endgroup
}
\makeatother

% 页脚右下方常态化水印(署名 + 仓库地址,加粗黑色),已在 main/plain 分页样式中定义,所有页面均生效.

\begin{document}
\mainmatter
\pagestyle{main}
\setcounter{chapter}{3}
\def\BookEnd{\end{document}}
\fi

\chapter{一元微分学与一元积分学}
\begin{ti}
【拓展1】$Rolle$定理的$Samuleson$证明
\end{ti}

\textbf{证明}:设$f\in C^0[a,b],\text{且在开区间}(a,b)\text{可导},f(a)=f(b),\text{则要证的就是}\;\exists\xi,\;\text{s.t.}\,f'(\xi)=0,$ \\

构造$F(x)=f(x)-f(x+\dfrac{b-a}{2}),F\in C[a,\dfrac{a+b}{2}],F(a)=-F(\dfrac{a+b}{2}),$ \\

$\therefore\exists\,a_1\in[a,\dfrac{a+b}{2}],\;\text{s.t.}\,F(a_1)=0,$ \\

记$b_1=a_1+\dfrac{b-a}{2},\,\,\text{则}f(a_1)=f(b_1),\,\,\text{且}b_1-a_1=\dfrac{b-a}{2},$ \\

以此类推地做下去,可以得到闭区间套$\{[a_n,b_n]\},\text{且}f(a_n)=f(b_n),b_n-a_n=\dfrac{b-a}{2^n},$ \\

由闭区间套原理,这些闭区间有唯一公共点$\xi,\text{且}\displaystyle\lim_{n\to\infty}a_n=\lim_{n\to\infty}b_n=\xi,a_n\leq\xi\leq b_n,$ \\

事实上,总可以使得$\displaystyle\xi\in(a,b)$.\\

若$\displaystyle\xi=a$,则由$\displaystyle a\leq a_n\leq\xi=a$可知$\displaystyle a_n=a$对每个$\displaystyle n$成立, \\

从而$\displaystyle f(b_1)=f(a)=f(b_2)$, \\

又$\displaystyle b_2=a+\dfrac{b-a}{4}<b_1=a+\dfrac{b-a}{2}$,故$\displaystyle[b_2,b_1]\subset(a,b)$, \\

以$\displaystyle[b_2,b_1]$为新的初始区间,对$\displaystyle f$重新进行上述构造, \\

所得闭区间套的公共点属于$\displaystyle[b_2,b_1]\subset(a,b)$, \\

若$\displaystyle\xi=b$,则同理由$\displaystyle\xi=b\leq b_n\leq b$可知$\displaystyle b_n=b$对每个$\displaystyle n$成立, \\

从而$\displaystyle f(a_1)=f(b)=f(a_2)$,且$\displaystyle[a_1,a_2]\subset(a,b)$, \\

以$\displaystyle[a_1,a_2]$为新的初始区间重新进行上述构造即可, \\

因此总能得到公共点$\displaystyle\xi\in(a,b)$的闭区间套, \\

下面证明一个引理:若有$x_n\leq c\leq y_n,\displaystyle\lim_{n\to\infty}x_n=\lim_{n\to\infty}y_n=c,f\text{在}c\text{处可导},$ \\
$\displaystyle\text{则}f'(c)=\lim_{n\to\infty}\dfrac{f(x_n)-f(y_n)}{x_n-y_n}.$ \\

不妨设$f'(c)=0,$否则令$g(x)=f(x)-f'(c)x,\;\text{则}g'(c)=0,$ \\

$\displaystyle\lim_{n\to\infty}\dfrac{f(x_n)-f(y_n)}{x_n-y_n}=\lim_{n\to\infty}\dfrac{g(x_n)-g(y_n)}{x_n-y_n}+f'(c),$ \\

$\because f'(c)=0,\quad\therefore\forall\,\varepsilon>0,\;\exists\,N,\;n>N\text{时},$ \\

$\left\lvert f(c)-f(x_n)\right\rvert<\varepsilon(c-x_n),\left\lvert f(y_n)-f(c)\right\rvert<\varepsilon(y_n-c),$ \\

$\therefore\left\lvert f(x_n)-f(y_n)\right\rvert\leq\left\lvert f(c)-f(x_n)\right\rvert+\left\lvert f(y_n)-f(c)\right\rvert<\varepsilon(y_n-x_n),$ \\

即$\forall\,\varepsilon>0,\;\exists\,N,\;n>N\text{时},\;\left\lvert \dfrac{f(x_n)-f(y_n)}{x_n-y_n}\right\rvert<\varepsilon,$ \\

$\therefore\displaystyle\lim_{n\to\infty}\dfrac{f(x_n)-f(y_n)}{x_n-y_n}=f'(c)=0,$得证, \\
 
由引理,$f'(\xi)=\displaystyle\lim_{n\to\infty}\dfrac{f(a_n)-f(b_n)}{a_n-b_n}=0$,命题得证. \\

\textbf{注}:以上证明并未用到最值定理与极值的思想,事实上,由$Samuleson$证明得到的$\xi$不一定是极值点, \\

如$f(x)=\begin{cases}
    x^3\sin{\dfrac{1}{x}}&\quad x\in[-\dfrac{1}{\pi},0)\cup(0,\dfrac{1}{\pi}] \\
    0&\quad x=0 \\
\end{cases}$

\begin{ti}
【拓展2】$Lagrange$定理的面积证法
\end{ti}

\textbf{证明}:令$S(x)=\dfrac{1}{2}\begin{vmatrix}
    1 &1 &1 \\
    a &b &x \\
    f(a) &f(b) &f(x)
\end{vmatrix},$ \\

则$S(x)=\dfrac{1}{2}\begin{vmatrix}
    1 &0 &0 \\
    a &b-a &x-a \\
    f(a) &f(b)-f(a) &f(x)-f(a)
\end{vmatrix}=\dfrac{1}{2}\begin{vmatrix}
     b-a &x-a \\
     f(b)-f(a) &f(x)-f(a)
\end{vmatrix},$ \\

$\therefore\,S(a)=S(b)=0,$由$Rolle$定理,$\exists\,\xi\in(a,b),\;\text{s.t.}\,S'(\xi)=0,$ \\

由行列式求导法则,$2S'(x)=\begin{vmatrix}
    0 &x-a \\
    0 &f(x)-f(a)
\end{vmatrix}+\begin{vmatrix}
    b-a &1 \\
    f(b)-f(a) &f'(x)
\end{vmatrix},$ \\

$\therefore\,2S'(\xi)=\begin{vmatrix}
    b-a &1 \\
    f(b)-f(a) &f'(\xi)
\end{vmatrix}=0\Rightarrow f'(\xi)=\dfrac{f(b)-f(a)}{b-a}.$ \\

\begin{minipage}[t]{0.59\textwidth}
\textbf{注}:这个证法来自于几何直观,我们要找函数上切线与给定连线平行的点 ,$S(x)$是$(a,f(a)),(b,f(b)),(x,f(x))$三点连成的三角形的有向面积,当$(x,f(x))$沿着平行于$(a,f(a)),(b,f(b))$的连线的直线移动时,$S(x)$不变,
猜想:$S(x)$的驻点处,切线与该连线平行
\end{minipage}\hfill
\begin{minipage}[t]{0.37\textwidth}
\begin{tikzpicture}[x=0.72cm,y=0.72cm,baseline=(current bounding box.north)]
    \draw[bookteal!55,line width=1pt] (0.35,0.35) -- (2.90,2.05) -- (4.35,0.75) -- cycle;
    \draw[bookteal!75,line width=0.8pt]
        (0.35,0.35) .. controls (0.95,1.35) and (2.30,1.99) .. (2.90,2.05)
        .. controls (3.12,2.072) and (3.25,2.08) .. (3.35,2.02)
        .. controls (3.72,1.80) and (4.08,1.02) .. (4.35,0.75);
    \draw[bookgray!80,dashed,line width=0.55pt]
        (1.80,1.94) -- (3.80,2.14);
    \fill[bookblue] (0.35,0.35) circle (1.4pt);
    \fill[bookblue] (2.90,2.05) circle (1.4pt);
    \fill[bookblue] (4.35,0.75) circle (1.4pt);
    \node[below left,font=\scriptsize] at (0.35,0.35) {$(a,f(a))$};
    \node[above left,font=\scriptsize] at (2.90,2.05) {$(x,f(x))$};
    \node[below right,font=\scriptsize] at (4.35,0.75) {$(b,f(b))$};
\end{tikzpicture}

\vspace{2pt}
{\centering\small\color{bookgray}图4-1\par}
\end{minipage}
\par

\begin{ti}
【题3】设$f\in C[0,1],\text{在}(0,1)\text{可微},f(0)=0,f(1)=1,\text{设}k_1,k_2,...,k_n\text{是满足}$ \\
$k_1+k_2+...+k_n=1$的$n$个正数,证明:在$(0,1)$中存在$n$个互不相同的数$t_1,t_2,...,t_n,\text{使得}\dfrac{k_1}{f'(t_1)}+...+\dfrac{k_n}{f'(t_n)}=1$
\end{ti}

\textbf{证明}:由介值定理,$\exists\,x_1\in(0,1),\;\text{s.t.}\,f(x_1)=k_1,$ \\

同理$\exists\,x_2\in(x_1,1),\;\text{s.t.}\,f(x_2)=k_1+k_2,$ \\

以此类推,可以在$[0,1]$中插入一系列数:$0=x_0<x_1<x_2<...<x_{n-1}<x_n=1,$ \\

同时$f(x_m)=\displaystyle\sum_{i=1}^{m}k_i$,在$[x_{i-1},x_i]\text{上,由}Lagrange$中值定理, \\

$\exists\,t_i\in(x_{i-1},x_i),\;\text{s.t.}\,f(x_i)-f(x_{i-1})=f'(t_i)(x_i-x_{i-1}),i=1,2,...,n,$ \\

$\therefore\displaystyle\sum_{i=1}^{n}\dfrac{k_i}{f'(t_i)}=\sum_{i=1}^{n}(x_i-x_{i-1})=1,\text{且}t_1,t_2,...,t_n$互不相同,得证.

\begin{ti}
【题4】设$f(x)=\displaystyle\sum_{k=1}^{n}c_ke^{\lambda_k x},\text{其中}\lambda_1,\lambda_2,...,\lambda_n$为互异实数,$c_1,c_2,...,c_n$不同时为$0$, \\

证明:$f$的零点个数小于$n$
\end{ti}

\textbf{证明}:$n=1\text{时},f(x)=c_1e^{\lambda_1 x}$无零点,符合要求, \\

假设$n=m$时结论成立,$n=m+1$时,令$g(x)=e^{-\lambda_{m+1} x}f(x)=\displaystyle\sum_{k=1}^{m+1}c_ke^{(\lambda_k-\lambda_{m+1})x},$ \\

则$g'(x)=\displaystyle\sum_{k=1}^{m+1}c_k(\lambda_k-\lambda_{m+1})e^{(\lambda_k-\lambda_{m+1})x}=\sum_{k=1}^{m}c_k(\lambda_k-\lambda_{m+1})e^{(\lambda_k-\lambda_{m+1})x},$ \\

记$d_k=c_k(\lambda_k-\lambda_{m+1}),\;\lambda_k'=\lambda_k-\lambda_{m+1},$ \\

那么$d_1,...,d_m$不同时为$0$,$\lambda_1',...,\lambda_m'$为互异实数,又$g'(x)=\displaystyle\sum_{k=1}^{m}c_k(\lambda_k-\lambda_{m+1})e^{(\lambda_k-\lambda_{m+1})x}\rlap{,}$ \\

由归纳假设可知,$g'(x)$零点个数小于$m$,\\

但如果$f$有大于等于$m+1$个零点,则$g$也有大于等于$m+1$个零点, \\

从而$g'$有至少$m$个零点,矛盾!$\quad\therefore f$的零点个数小于$m+1,$ \\

由数学归纳法可知,$f$的零点个数小于$n.$

\begin{ti}
【题5】$\;f\text{在}[a,b]\text{上连续},(a,b)\text{上二阶可导},f(a)=f(b)=0,\int_{a}^{b}f(x)\,dx=0,$ \\
证明:$\exists\,\zeta\in[a,b],\;\text{s.t.}f(\zeta)=f''(\zeta)$
\end{ti}

\textbf{证明}:若$f(x)_{min}=0,\text{则由}\int_{a}^{b}f(x)\,dx=0\Rightarrow f\equiv 0,$ \\

则$\forall\zeta\in(a,b),\text{都有}f(\zeta)=f''(\zeta)$,同理$f(x)_{max}=0$时也有$f(\zeta)=f''(\zeta),$ \\

下设$f(x_1)=f(x)_{min}<0,f(x_2)=f(x)_{max}>0,x_1,x_2\in(a,b),$ \\

由零值定理,$\exists\,x_0\in(a,b),\;\text{s.t.}f(x_0)=0,$ \\

令$G(x)=e^xf(x),\text{则}G(a)=G(x_0)=G(b)=0,G'(x)=e^x[f(x)+f'(x)],$ \\

由$Rolle$定理,$\;\exists\,\xi_1\in(a,x_0),\xi_2\in(x_0,b),\;\text{s.t.}\,G'(\xi_1)=G'(\xi_2)=0,$ \\

令$F(x)=e^{-x}[f(x)+f'(x)],\text{则}F(\xi_1)=F(\xi_2)=0,$ \\

$\because f\text{在}[a,b]\text{上连续},\text{在}(a,b)\text{上二阶可导},\quad\therefore F\text{在}[\xi_1,\xi_2]\text{上连续},\text{在}(\xi_1,\xi_2)\text{上可导},$ \\

由$Rolle$定理,$\;\exists\,\zeta\in(\xi_1,\xi_2),\;\text{s.t.}\,F'(\zeta)=0$, \\

又$F'(x)=e^{-x}[f''(x)-f(x)],\quad\therefore f''(\zeta)=f(\zeta)$, \\

综上,$\exists\,\zeta\in[a,b],\;\text{s.t.}f(\zeta)=f''(\zeta).$ 

\begin{ti}
【题6】设$f\text{在}[a,b]\text{上可微,在}(a,b)\text{上二阶可微},$ 证明:$\exists\,\xi\in(a,b),\;\text{s.t.}\,f'(b)-f'(a)=f''(\xi)(b-a)$
\end{ti}

\textbf{证明}:假设不存在这样的$\xi$,由$Darboux$定理, \\

$f''(x)>\dfrac{f'(b)-f'(a)}{b-a}\text{恒成立或}f''(x)<\dfrac{f'(b)-f'(a)}{b-a}\text{恒成立}$, \\

不妨设$f''(x)>\dfrac{f'(b)-f'(a)}{b-a},\;\forall\,x\in(a,b)$, \\

设$\lambda=\dfrac{f'(b)-f'(a)}{b-a},\;g(x)=f'(x)-\lambda x,\text{则}x\in(a,b)\text{时},g'(x)=f''(x)-\lambda>0$, \\

$\therefore\,g(x)\text{在}(a,b)$上严格单调递增,且$g(a)=g(b)$, \\

$\therefore\,g(a^+),g(b^-)$均存在,即$\displaystyle\lim_{x\to a^+}f'(x),\lim_{x\to b^-}f'(x)$均存在, \\

下面证明一个重要引理:若$\displaystyle\lim_{x\to a^+}f'(x)$存在,则必然等于$f'(a)$, \\

事实上,由$Lagrange$中值定理,$\exists\,\xi_x\in(a,x),\;\text{s.t.}\,\dfrac{f(x)-f(a)}{x-a}=f'(\xi_x)$, \\

$x\rightarrow a^+\text{时},\xi_x\rightarrow a^+$,对上式两边取极限即得到上述引理. \\

由上面的引理,可知$g$在$a,b$处均连续,即$g\in C[a,b]$, \\

$\therefore g\text{在}[a,b]$上严格单调递增,这与$g(a)=g(b)$矛盾! \\

综上,$\exists\,\xi\in(a,b),\;\text{s.t.}\,f'(b)-f'(a)=f''(\xi)(b-a)$. \\

\textbf{注意}:本题并未要求$f'\in C[a,b]$,本题的意义在于,对于导函数来说,$Lagrange$定理中在闭区间上连续的要求不是必要的.

\begin{ti}
【题7】设$f\in C^1[a,b],\int_{a}^{b}f(x)\,dx=\int_{a}^{b}xf(x)\,dx=0,$求证:$\exists\,x_1,x_2\in(a,b),x_1\neq x_2,\;\text{s.t.}\,f(x_1)=f(x_2)=0$
\end{ti}

\textbf{证明}:若$f(x)\text{在}(a,b)$上无零点,由介值定理,$f(x)>0\text{或}f(x)<0$恒成立, \\

这与$\displaystyle\int_{a}^{b}f(x)\,dx=0$矛盾,$\quad\therefore\exists\,c\in(a,b),\;\text{s.t.}\,f(c)=0$, \\

若$f(x)$在$(a,b)$中只有这一个零点, \\

如果$f(x)$在$c$两边同号,则与$\displaystyle\int_{a}^{b}f(x)\,dx=0$矛盾, \\

$\therefore\,f(x)\text{在}c$两边异号, \\

不妨设$x\in(a,c)\text{时},f(x)<0,\;x\in(c,b)\text{时},f(x)>0$, \\

则$(x-c)f(x)>0$对$\forall\,x\in(a,b)$恒成立,$\;\therefore\displaystyle\int_{a}^{b}(x-c)f(x)\,dx>0$, \\

而$\displaystyle\int_{a}^{b}(x-c)f(x)\,dx=\int_{a}^{b}xf(x)\,dx-c\int_{a}^{b}f(x)\,dx=0$,矛盾! \\

所以$f$至少两个零点,即:$\exists\,x_1,x_2\in(a,b),x_1\neq x_2,\;\text{s.t.}\,f(x_1)=f(x_2)=0$.

\begin{ti}
【题8】设$f$在$[a,b]$上二阶可微,$f'(a)=f'(b)=0$,证明:$\exists\,\xi\in(a,b),\;\text{s.t.}\,\left\lvert f''(\xi)\right\rvert\geq\dfrac{4}{(b-a)^2}\left\lvert f(b)-f(a)\right\rvert $
\end{ti}

\textbf{证明}:将$f(x)$分别在$a$处和$b$处展开, \\

$f(x)=f(a)+f'(a)(x-a)+\dfrac{f''(\xi_1)}{2}(x-a)^2=f(a)+\dfrac{f''(\xi_1)}{2}(x-a)^2$, \\

$f(x)=f(b)+f'(b)(x-b)+\dfrac{f''(\xi_2)}{2}(x-b)^2=f(b)+\dfrac{f''(\xi_2)}{2}(x-b)^2$, \\

代入$x=\dfrac{a+b}{2},f(\dfrac{a+b}{2})=f(a)+\dfrac{f''(\xi_1)}{2}\cdot\dfrac{(b-a)^2}{4}=f(b)+\dfrac{f''(\xi_2)}{2}\cdot\dfrac{(b-a)^2}{4}$, \\

则$\left\lvert f(b)-f(a)\right\rvert\cdot\dfrac{4}{(b-a)^2}=\dfrac{1}{2}\left\lvert f''(\xi_2)-f''(\xi_1)\right\rvert$ $\leq\dfrac{1}{2}\left(\left\lvert f''(\xi_1)\right\rvert+\left\lvert f''(\xi_2)\right\rvert\right)$, \\

取$\xi\in\{\xi_1,\xi_2\},\text{使得}\left\lvert f''(\xi)\right\rvert=\max\{\left\lvert f''(\xi_1)\right\rvert,\left\lvert f''(\xi_2) \right\rvert\}$, \\

则$\left\lvert f''(\xi)\right\rvert\geq\dfrac{4}{(b-a)^2}\left\lvert f(b)-f(a)\right\rvert$,得证.

\begin{ti}
【题9】设$f$在$(a,b)$上可微,且在某点$\xi\in(a,b)$处有$f''(\xi)>0$,证明:$\exists\,x_1,x_2\in(a,b),\;\text{s.t.}\,\dfrac{f(x_2)-f(x_1)}{x_2-x_1}=f'(\xi)$
\end{ti}

\textbf{证明}:令$F(x)=f(x)-f'(\xi)x,\;\text{则}F'(\xi)=f'(\xi)-f'(\xi)=0,F''(\xi)=f''(\xi)>0,$ \\

由二阶导数的定义,$\therefore\exists\,\delta,\;\text{s.t.}\,x\in[\xi-\delta,\xi)\text{时},F'(x)<0;\;x\in(\xi,\xi+\delta]\text{时},F'(x)>0,$ \\

$\therefore\,x\in[\xi-\delta,\xi)\text{时},F(x)\text{严格单调递减},x\in(\xi,\xi+\delta]\text{时},F(x)\text{严格单调递增},$ \\

设$m=\min\{F(\xi-\delta),F(\xi+\delta)\},\text{则}F(\xi)<m,$ \\

取$\eta\in(F(\xi),m),\text{则}\exists\,x_1\in(\xi-\delta,\xi),x_2\in(\xi,\xi+\delta),\;\text{s.t.}\,F(x_1)=F(x_2)=\eta,$ \\

$\therefore\,0=\dfrac{F(x_2)-F(x_1)}{x_2-x_1}=\dfrac{f(x_2)-f(x_1)}{x_2-x_1}-f'(\xi),$ \\

$\therefore\,\dfrac{f(x_2)-f(x_1)}{x_2-x_1}=f'(\xi)$,这样的$x_1,x_2$即为所求. \\

\textbf{注}:本题虽然简单,但是可以与$Lagrange$中值定理对比来看,事实上,如果没有其他条件,那么任意的$\xi\in(a,b)$,不一定存在$x_1,x_2$,使得$\dfrac{f(x_2)-f(x_1)}{x_2-x_1}=f'(\xi)$,
读者可以自行思考存在$x_1,x_2$使得$\dfrac{f(x_2)-f(x_1)}{x_2-x_1}=f'(\xi)$的充要条件是什么.

\begin{ti}
【题10】设$f$在$[a,b]$上可微,$f'(a)=f'(b)$,证明:$\exists\,\xi\in(a,b),\;\text{s.t.}\,f'(\xi)=\dfrac{f(\xi)-f(a)}{\xi-a}$
\end{ti}

\textbf{证明}:令$F(x)=f(x)-f(a)-f'(a)(x-a),\text{则}F(a)=0,F'(a)=F'(b)=0,$ \\

\begin{minipage}[t]{0.58\textwidth}
考虑本题的几何含义, \\

定义$H(x)=\begin{cases}
    \dfrac{F(x)}{x-a}&\quad (a<x\leq b) \\
    0&\quad (x=a) \\
\end{cases}$ \\[3pt]

易得$H(a^+)=H(a)$, \\

$\therefore H\text{在}[a,b]\text{上连续}$, $H\text{在}(a,b)\text{上可导}$, \\
\end{minipage}\hfill
\begin{minipage}[t]{0.38\textwidth}
\centering
\resizebox{0.96\linewidth}{!}{%
\begin{tikzpicture}[x=0.90cm,y=1.05cm,baseline=(current bounding box.north)]
    \draw[bookteal!80,line width=0.9pt]
        (0.35,0.45) .. controls (0.85,1.55) and (1.45,2.25) .. (2.15,1.75)
        .. controls (2.75,1.30) and (3.80,1.5253) .. (4.20,1.65)
        .. controls (4.60,1.7747) and (5.15,1.82) .. (5.70,1.55);
    \draw[bookgray!85,line width=0.8pt]
        (-0.10,0.3097) -- (5.70,2.1175);
    \fill[bookblue] (0.35,0.45) circle (1.5pt);
    \fill[bookblue] (4.20,1.65) circle (1.5pt);
    \node[below left,font=\scriptsize] at (0.35,0.45) {$(a,f(a))$};
    \node[below right,font=\scriptsize] at (4.20,1.65) {$(x,f(x))$};
\end{tikzpicture}%
}

\vspace{2pt}
{\centering\small\color{bookgray}图4-2\par}
\end{minipage}
\par


且$x\in(a,b]\text{时},\;H'(x)=\dfrac{(x-a)f'(x)-f(x)+f(a)}{(x-a)^2},$ \\

$\therefore\,H(b)=\dfrac{F(b)}{b-a},\;H'(b)=\dfrac{(b-a)f'(b)-f(b)+f(a)}{(b-a)^2}=-\dfrac{F(b)}{(b-a)^2}=-\dfrac{H(b)}{b-a}.$ \\

\romannum{1} $\quad H(b)>0,\text{则}H'(b)<0,\text{则}\exists\,c\in(a,b),\;\text{s.t.}\,H(c)>H(b),$ \\

又$H(c)>H(b)>0=H(a),\;\therefore\,H$在区间内部取到最大值,\\

即$\exists\,\xi\in(a,b),\;\text{s.t.}\,H'(\xi)=0.$ \\

\romannum{2} $\quad H(b)<0,\text{则}H'(b)>0,\text{则}\exists\,c\in(a,b),\;\text{s.t.}\,H(c)<H(b),$ \\

又$H(c)<H(b)<0=H(a),\;\therefore\,H$在区间内部取到最小值, \\

即$\exists\,\xi\in(a,b),\;\text{s.t.}\,H'(\xi)=0.$ \\

\romannum{3} $\quad H(b)=0$,则由$Rolle$定理,$\exists\,\xi\in(a,b),\;\text{s.t.}\,H'(\xi)=0.$ \\

综上,$\exists\,\xi\in(a,b),\;\text{s.t.}\,H'(\xi)=0$,即$\exists\,\xi\in(a,b),\;\text{s.t.}\,f'(\xi)=\dfrac{f(\xi)-f(a)}{\xi-a}.$ 

\begin{ti}
【题11】设$f$在$[0,+\infty)$上可微,$f(0)=0,\text{且}\exists\,c>0,\;\text{s.t.}\,\left\lvert f'(x)\right\rvert\leq c\left\lvert f(x)\right\rvert,\forall\,x\in[0,+\infty),\;$
证明:$f(x)\equiv 0$
\end{ti}

\textbf{证明}:易得$\left\lvert f\right\rvert\in C[0,\dfrac{1}{2c}]$,由连续函数最值定理,设$\left\lvert f(x_0)\right\rvert=\displaystyle\max_{x\in[0,\frac{1}{2c}]}\left\lvert f(x)\right\rvert$, \\

$\therefore\left\lvert f(x_0)\right\rvert=\left\lvert f(x_0)-f(0)\right\rvert=\left\lvert f'(\xi)\right\rvert x_0\leq c\left\lvert f(\xi)\right\rvert x_0$ \\

$\leq\dfrac{1}{2}\left\lvert f(\xi)\right\rvert\leq\dfrac{1}{2}\left\lvert f(x_0)\right\rvert\Longrightarrow \left\lvert f(x_0)\right\rvert=0,$ $\;\therefore\,f(x)\equiv 0,\;\forall\,x\in[0,\dfrac{1}{2c}]$, \\

假设$x\in[\dfrac{k-1}{2c},\dfrac{k}{2c}]\text{时},f(x)\equiv 0,\text{记}\left\lvert f(x_1)\right\rvert=\displaystyle\max_{x\in[\frac{k}{2c},\frac{k+1}{2c}]}\left\lvert f(x)\right\rvert,$ \\

$\therefore\left\lvert f(x_1)\right\rvert=\left\lvert f(x_1)-f(\dfrac{k}{2c})\right\rvert=\left\lvert f'(\xi_1)\right\rvert(x_1-\dfrac{k}{2c})\leq c\left\lvert f(\xi_1)\right\rvert\cdot\dfrac{1}{2c}$ \\

$\leq\dfrac{1}{2}\left\lvert f(x_1)\right\rvert\Longrightarrow\left\lvert f(x_1)\right\rvert=0,\;$ 即$f(x)\equiv 0,\;\forall\,x\in[\dfrac{k}{2c},\dfrac{k+1}{2c}],$ \\

由数学归纳法,$f(x)\equiv 0,\;\forall\,x\in[\dfrac{n}{2c},\dfrac{n+1}{2c}],\;\text{对}\forall\,n\in\mathbb{N}\text{恒成立}$, \\

即$f(x)\equiv 0,\;\forall\,x\in\displaystyle\bigcup_{n=0}^{\infty}[\dfrac{n}{2c},\dfrac{n+1}{2c}]$,\\

即$f(x)\equiv 0,\;\forall\,x\in[0,+\infty)$,证毕. \\

\textbf{另解}:由题意可得$f(x)=\displaystyle\int_{0}^{x}f'(t)\,dt,$ \\

$\displaystyle\therefore\left\lvert f(x)\right\rvert=\left\lvert \int_{0}^{x}f'(t)\,dt\right\rvert\leq\int_{0}^{x}\left\lvert f'(t)\right\rvert \,dt\leq c\int_{0}^{x}\left\lvert f(t)\right\rvert\,dt $, \\

$\forall\,x_0\in(0,+\infty),\;f\in C[0,x_0],\text{记}M=\displaystyle\max_{t\in[0,x_0]}\left\lvert f(t)\right\rvert,$ \\

则$\displaystyle\left\lvert f(x)\right\rvert\leq c\int_{0}^{x}\left\lvert f(t)\right\rvert \,dt\leq c\cdot Mx,\;\forall\,x\in[0,x_0],$ \\

再将此式带回,可得$\displaystyle\left\lvert f(x)\right\rvert\leq c\int_{0}^{x}\left\lvert f(t)\right\rvert \,dt\leq c^2\int_{0}^{x}Mt\,dt=\dfrac{Mc^2}{2}x^2,$ \\

以此归纳地做下去,可得$x\in[0,x_0]\text{时},\left\lvert f(x)\right\rvert\leq\dfrac{Mc^nx^n}{n!},\text{对}\forall\,n\in\mathbb{N}^+$恒成立, \\

$\therefore\,x\in[0,x_0]\text{时},\left\lvert f(x)\right\rvert\leq\displaystyle\lim_{n\to\infty}\dfrac{Mc^nx^n}{n!}=0,\text{即}f(x)=0,$ \\

由$x_0$的任意性可知,$f\equiv 0$,证毕.



\BookEnd
